% ============================================================
\chapter*{Preface}
\addcontentsline{toc}{chapter}{Preface}
% ============================================================

\section*{Course Objectives}

Databases are one of the fundamental pillars of modern computer science.
From the early file management systems of the 1960s to the explosion of
NoSQL databases and data lakes, the ability to store, organize, and query
data reliably and efficiently remains a central challenge for every
software application.

This course is designed for second-year undergraduate students (L2 / sophomore level)
in computer science, mathematics-computer science, or data science. It assumes
basic familiarity with programming (Python) and discrete mathematics (sets,
relations, propositional logic).

Upon completion of this course, the student will be able to:
\begin{itemize}
    \item Design a conceptual schema (Entity-Relationship model) and translate
          it into a normalized relational schema.
    \item Write complex SQL queries involving joins, aggregations, subqueries,
          and window functions.
    \item Understand and apply normal forms (1NF through BCNF) to eliminate
          update anomalies.
    \item Explain the ACID properties and concurrency control mechanisms.
    \item Choose and create appropriate indexes to optimize performance.
    \item Design a simple ETL pipeline using Python and SQL.
    \item Compare NoSQL models (document, key-value, column, graph)
          with the relational model.
\end{itemize}

\section*{Course Organization}

The course is structured into eleven chapters, progressing from theoretical
foundations to practical applications:

\begin{enumerate}
    \item \textbf{Introduction to Relational Databases} ---
          History, core concepts, DBMS architecture.
    \item \textbf{Entity-Relationship Model} ---
          Conceptual schema design, ER diagrams, mapping to relational schemas.
    \item \textbf{Relational Algebra} ---
          Formal operators: selection, projection, join, division.
    \item \textbf{SQL --- Basic Queries} ---
          \texttt{SELECT}, \texttt{FROM}, \texttt{WHERE}, sorting and filtering.
    \item \textbf{SQL --- Joins and Subqueries} ---
          Inner, outer, cross joins; correlated subqueries.
    \item \textbf{SQL --- Aggregation and Window Functions} ---
          \texttt{GROUP BY}, \texttt{HAVING}, \texttt{ROW\_NUMBER}, \texttt{RANK}.
    \item \textbf{Normalization} ---
          Functional dependencies, 1NF through BCNF, lossless decomposition.
    \item \textbf{Transactions, Integrity, Constraints} ---
          ACID properties, isolation levels, locking.
    \item \textbf{Indexing and Query Optimization} ---
          B-trees, hash indexes, query execution plans.
    \item \textbf{Data Pipelines and ETL} ---
          Extract, Transform, Load; Python/SQL integration.
    \item \textbf{NoSQL Databases} ---
          Document, key-value, column, graph; the CAP theorem.
\end{enumerate}

\section*{Pedagogical Approach}

Each chapter follows a consistent structure:

\begin{itemize}
    \item \textbf{Theory} --- Formal definitions, theorems, and propositions
          with proofs or justifications.
    \item \textbf{Detailed Examples} --- Every concept is illustrated with one
          or more concrete examples, often using a ``University'' database
          as a running example.
    \item \textbf{Executable SQL Code} --- Queries are presented with their
          expected output, testable on SQLite, PostgreSQL, or MySQL.
    \item \textbf{Python Integration} --- \texttt{sqlite3} and \texttt{pandas}
          scripts demonstrate programmatic usage.
    \item \textbf{Exercises} --- Exercises of increasing difficulty, from
          comprehension checks to open-ended problem solving.
\end{itemize}

\section*{Running Example Database}

Throughout this course, we use a \textbf{University} database with the
following tables:

\begin{codeblock}[title=\textbf{University Database Schema}]
\begin{minted}{sql}
CREATE TABLE Students (
    student_id   INTEGER PRIMARY KEY,
    last_name    TEXT NOT NULL,
    first_name   TEXT NOT NULL,
    birth_date   DATE,
    major        TEXT
);

CREATE TABLE Courses (
    course_id    INTEGER PRIMARY KEY,
    title        TEXT NOT NULL,
    credits      INTEGER DEFAULT 3,
    prof_id      INTEGER REFERENCES Professors(prof_id)
);

CREATE TABLE Professors (
    prof_id      INTEGER PRIMARY KEY,
    last_name    TEXT NOT NULL,
    first_name   TEXT NOT NULL,
    department   TEXT
);

CREATE TABLE Enrollments (
    student_id   INTEGER REFERENCES Students(student_id),
    course_id    INTEGER REFERENCES Courses(course_id),
    grade        REAL,
    year         INTEGER,
    PRIMARY KEY (student_id, course_id)
);
\end{minted}
\end{codeblock}

This schema will be enriched throughout the chapters to illustrate advanced
concepts (inheritance, ternary relationships, historization).

\section*{Sample Data}

\begin{codeblock}[title=\textbf{Inserting Sample Data}]
\begin{minted}{sql}
INSERT INTO Professors VALUES (1, 'Smith', 'Mary', 'Computer Science');
INSERT INTO Professors VALUES (2, 'Johnson', 'John', 'Mathematics');
INSERT INTO Professors VALUES (3, 'Williams', 'Sophie', 'Computer Science');

INSERT INTO Courses VALUES (101, 'Databases', 4, 1);
INSERT INTO Courses VALUES (102, 'Algorithms', 3, 2);
INSERT INTO Courses VALUES (103, 'Networks', 3, 3);
INSERT INTO Courses VALUES (104, 'Statistics', 3, 2);

INSERT INTO Students VALUES (1, 'Brown', 'Alice', '2004-03-15', 'CS');
INSERT INTO Students VALUES (2, 'Davis', 'Bob', '2003-11-22', 'CS');
INSERT INTO Students VALUES (3, 'Miller', 'Claire', '2004-07-08', 'Math');
INSERT INTO Students VALUES (4, 'Wilson', 'David', '2003-01-30', 'CS');
INSERT INTO Students VALUES (5, 'Taylor', 'Emma', '2004-09-12', 'Math');

INSERT INTO Enrollments VALUES (1, 101, 88.5, 2025);
INSERT INTO Enrollments VALUES (1, 102, 76.0, 2025);
INSERT INTO Enrollments VALUES (2, 101, 72.0, 2025);
INSERT INTO Enrollments VALUES (2, 103, 91.0, 2025);
INSERT INTO Enrollments VALUES (3, 102, 95.5, 2025);
INSERT INTO Enrollments VALUES (3, 104, 82.0, 2025);
INSERT INTO Enrollments VALUES (4, 101, 55.0, 2025);
INSERT INTO Enrollments VALUES (5, 104, 97.0, 2025);
\end{minted}
\end{codeblock}

\section*{Recommended Tools}

\begin{itemize}
    \item \textbf{SQLite} --- Lightweight, serverless, ideal for learning.
          Available via the \texttt{sqlite3} command line or the Python
          \texttt{sqlite3} module.
    \item \textbf{PostgreSQL} --- Full-featured DBMS, SQL-standard compliant,
          recommended for advanced exercises (window functions, concurrent
          transactions).
    \item \textbf{DB Browser for SQLite} --- GUI for visualizing and
          manipulating SQLite databases.
    \item \textbf{DBeaver} --- Universal client supporting PostgreSQL,
          MySQL, SQLite, and others.
    \item \textbf{Python 3.x} --- With modules \texttt{sqlite3} (standard library),
          \texttt{psycopg2} (PostgreSQL), \texttt{pandas} (data analysis).
\end{itemize}

\section*{Typographic Conventions}

\begin{itemize}
    \item SQL keywords are written in \texttt{UPPERCASE}:
          \texttt{SELECT}, \texttt{WHERE}, \texttt{JOIN}.
    \item Table and column names use \texttt{PascalCase} or
          \texttt{snake\_case} depending on context.
    \item Technical terms appearing for the first time are in \emph{italics}.
    \item Colored boxes indicate:
          \begin{itemize}
              \item important \textbf{definitions} and \textbf{theorems},
              \item \textbf{best practices} to remember,
              \item \textbf{anti-patterns} to avoid,
              \item \textbf{warnings} about common pitfalls.
          \end{itemize}
\end{itemize}

\section*{Prerequisites}

\begin{itemize}
    \item \textbf{Programming} --- Ability to write simple Python programs
          (variables, loops, functions, file I/O).
    \item \textbf{Mathematics} --- Basic set theory, relations, Cartesian
          products, propositional logic
          ($\land$, $\lor$, $\neg$, $\Rightarrow$).
    \item \textbf{Systems} --- Familiarity with a command-line terminal.
\end{itemize}

\section*{Supplementary Resources}

\begin{itemize}
    \item A.~Silberschatz, H.~Korth, S.~Sudarshan, \emph{Database System Concepts},
          McGraw-Hill, 7th edition.
    \item R.~Elmasri, S.~Navathe, \emph{Fundamentals of Database Systems},
          Pearson, 7th edition.
    \item R.~Ramakrishnan, J.~Gehrke, \emph{Database Management Systems},
          McGraw-Hill, 3rd edition.
    \item Official PostgreSQL Documentation: \url{https://www.postgresql.org/docs/}
    \item SQLite Documentation: \url{https://www.sqlite.org/docs.html}
    \item W3Schools SQL Tutorial: \url{https://www.w3schools.com/sql/}
\end{itemize}

\vfill
\begin{flushright}
\textit{Happy reading and happy learning!}
\end{flushright}
